% === §4.2 Objetivos específicos ===
% Redactado por el Investigador — Fase 3
% Mapeo 1:1 con los 3 subproblemas de §2.1
% Estructura: OE1 (Fundamentos y Caracterización de Datos) ↔ SP1
%            OE2 (Novedad Teórica y Metodológica en IA) ↔ SP2
%            OE3 (Transferencia Tecnológica, Despliegue y Validación) ↔ SP3
% Fases del alcance (§3): F1↔OE1, F2↔OE2, F3↔OE3 (validación), F4↔OE3 (transferencia)

Los tres objetivos específicos que se enuncian a continuación se derivan lógicamente del objetivo general y
mantienen una correspondencia biunívoca con los subproblemas identificados en la
sección~\ref{sec:problematica}, conformando una cadena de valor metodológica que conduce desde la
caracterización del aprendiz hasta la transferencia tecnológica con TRL~6/7.

\begin{enumerate}[leftmargin=*, itemsep=12pt]

    % --- OE1: Fundamentos y Caracterización de Datos (resuelve SP1) ---
    \item \textbf{Caracterizar y modelar computacionalmente el nivel de aprendizaje profundo
    —motivación intrínseca, argumentación técnica y resolución de problemas auténticos— en
    estudiantes de ingeniería que interactúan con herramientas de inteligencia artificial}
    \hfill \textbf{[SP1, F1]}

    \textit{Subproblema que resuelve:} SP1 — Diagnóstico y modelado multidimensional del
    aprendiz. La ausencia de instrumentos y modelos computacionales que permitan diagnosticar de
    manera sistemática, trazable y multidimensional el aprendizaje profundo será abordada mediante
    la construcción de un instrumento de diagnóstico validado psicométricamente y el desarrollo de
    un modelo computacional del aprendiz basado en técnicas de procesamiento de lenguaje natural,
    analítica del aprendizaje y aprendizaje automático supervisado, alimentado por un
    \emph{dataset} multimodal curado —trazas de interacción con código, textos argumentativos,
    soluciones a problemas abiertos de ingeniería y registros de simulación— recolectado en
    asignaturas reales de los programas objetivo.

    \textit{Forma de validación cuantitativa:} índices de consistencia interna del instrumento
    (alfa de Cronbach $\geq 0.80$ y omega de McDonald), validación de constructo mediante juicio
    de expertos (índice de validez de contenido, IVC $\geq 0.80$), y métricas de desempeño del
    modelo del aprendiz en la clasificación de niveles de aprendizaje profundo (F1-score macro,
    AUC-ROC, matriz de confusión con validación cruzada estratificada $k$-fold, $k=5$).
    \textit{Validación cualitativa:} análisis de contenido temático de producciones textuales y de
    código de los estudiantes, triangulado con observaciones de aula y diarios de campo del equipo
    investigador.

    % --- OE2: Novedad Teórica y Metodológica en IA (resuelve SP2) ---
    \item \textbf{Desarrollar un ecosistema de agentes autónomos basados en inteligencia
    artificial, pedagógicamente fundado y técnicamente validado, que articule funciones de tutoría
    personalizada, retroalimentación adaptativa, generación de problemas auténticos, asistencia en
    programación, simulación de sistemas ingenieriles y evaluación formativa}
    \hfill \textbf{[SP2, F2]}

    \textit{Subproblema que resuelve:} SP2 — Arquitectura agéntica autónoma para personalización
    del aprendizaje. La carencia de un ecosistema de agentes autónomos que orqueste de forma
    coherente las funciones didácticas requeridas para el aprendizaje profundo en ingeniería será
    superada mediante el diseño e implementación de una arquitectura multiagente con roles
    pedagógicamente fundados —agente tutor, agente evaluador, agente generador de problemas,
    agente simulador de sistemas eléctricos/electrónicos y agente asistente de programación—,
    orquestada bajo el Model Context Protocol (MCP) e integrada con grandes modelos de lenguaje
    (LLMs) preentrenados, mecanismos de memoria compartida y contexto pedagógico, y principios de
    \emph{scaffolding} adaptativo graduado según el nivel de competencia diagnosticado por el
    modelo del aprendiz (OE1).

    \textit{Forma de validación cuantitativa:} métricas de desempeño agéntico —tasa de
    finalización exitosa de tareas multi-paso en \emph{benchmarks} de tutoría en ingeniería,
    precisión de la retroalimentación generada evaluada contra rúbricas por pares expertos
    (acuerdo interevaluador $\kappa$ de Cohen $\geq 0.70$), latencia de respuesta extremo a
    extremo ($p_{95}$ en segundos), y tasa de alucinación factual en dominios ingenieriles
    (verificación automatizada contra fuentes canónicas)—, complementadas con pruebas de
    integración continua y cobertura funcional de los requisitos pedagógicos definidos en la fase
    de fundamentación. \textit{Validación cualitativa:} juicio de expertos en educación en
    ingeniería sobre la pertinencia pedagógica de la arquitectura agéntica, aplicando una rúbrica
    de alineación constructiva entre roles agénticos, estrategias didácticas y dimensiones del
    aprendizaje profundo.

    % --- OE3: Transferencia Tecnológica, Despliegue y Validación (resuelve SP3, cubre F3 y F4) ---
    \item \textbf{Validar empíricamente el impacto causal del ecosistema agéntico en el
    aprendizaje profundo de estudiantes de ingeniería mediante un diseño cuasiexperimental en
    entornos reales de aula universitaria, y desplegar el ecosistema en infraestructura híbrida
    alcanzando TRL~6/7 transferible al sector productivo colombiano}
    \hfill \textbf{[SP3, F3--F4]}

    \textit{Subproblema que resuelve:} SP3 — Validación empírica y transferencia tecnológica
    TRL~6/7. La falta de evidencia cuantitativa y cualitativa sobre el impacto causal del
    ecosistema agéntico, así como de una ruta de despliegue tecnológico transferible, será
    abordada en dos fases encadenadas. En la Fase~3 (validación), se ejecutará un diseño
    cuasiexperimental pre--post con grupo de control no equivalente en al menos dos asignaturas
    reales de pregrado y posgrado de los programas de Ingeniería Eléctrica, Electrónica, Ciencias
    de la Computación e Ingeniería de Software de la UNAL, recolectando datos longitudinales de
    aprendizaje profundo, interacción con el ecosistema y percepción de estudiantes y docentes.
    En la Fase~4 (transferencia), se desplegará el ecosistema en infraestructura híbrida
    —estaciones Apple Silicon locales para inferencia \emph{on-device} y VPS en nube para
    orquestación y analítica—, se documentarán la interoperabilidad (API, MCP, estándares LTI
    cuando sea factible) y los procedimientos de instalación y operación, se obtendrá el registro
    de software ante la Dirección Nacional de Derechos de Autor de Colombia, se publicará el
    código fuente bajo licencia abierta permisiva, y se ejecutará un plan de adopción con al
    menos un actor del sector productivo (empresa de software o centro de formación técnica del
    eje cafetero) que demuestre la operación del ecosistema en un entorno relevante.

    \textit{Forma de validación cuantitativa:} diferencia estadísticamente significativa
    ($p<0.05$) en las puntuaciones de rúbricas validadas de aprendizaje profundo entre los grupos
    experimental y control, con tamaño del efecto reportado ($d$ de Cohen, $\eta^2$ parcial);
    indicadores de adopción tecnológica (usuarios activos, sesiones completadas, tasa de
    retención); y verificación del nivel TRL~6 mediante lista de cotejo estandarizada, con
    proyección documentada a TRL~7 (demostración operativa con actor externo). \textit{Validación
    cualitativa:} análisis temático de entrevistas semiestructuradas a estudiantes y docentes,
    análisis de la calidad de la argumentación técnica en tareas auténticas de ingeniería mediante
    rúbricas de razonamiento científico, y grupos focales con actores del sector productivo para
    evaluar la pertinencia y usabilidad del ecosistema en contextos de formación técnica y
    empresarial.

\end{enumerate}
